Un four de céramiste dont l'intérieur est à \qty{\v{T}}{\degreeCelsius} possède un petit regard d'observation de \qty{\v{S}}{cm^2}. L'ouverture d'une cavité se comporte comme un corps noir. On néglige le rayonnement qui entre par l'ouverture.

\begin{enumerate}[label=\alph*)]
    \item Pourquoi l'ouverture d'une cavité se comporte-t-elle comme un corps noir ?
    \item Quelle puissance s'échappe par l'ouverture ?
    \item Quelle énergie est perdue de cette façon pendant une cuisson de \qty{\v{t}}{h} ? À quelle longueur d'onde le rayonnement qui s'échappe est-il maximal ?
\end{enumerate}
